% LAB BOREDOM — module L1, The Three Literatures. GUIDED DRAFTING (desktop recipe), Section 2 run.
% Session 2026-08-06. Author rulings: pack v2 is the bank; r1 ¶1 carried as working base.
% Pack: lab-boredom-L1-pack-v2 (+ -v2-quotes.json)  |  Spec: outline v7 §L1  |  Rulings: register v8.
% Protocol: Stage-3 guided batches (1–3 ¶¶, approval gate each batch); «Qnn» markers, mechanical
%   substitution; rendered copy = L1-DRAFT-v1-rendered.tex. MASTER FILE — edit this, re-render after.
%
% STANDING RULINGS ADDED THIS SESSION (author notes on batch 1):
%   (i) voice must match Part III Section 1 (desktop register, its cadence);
%   (ii) FIRST + LAST NAME at an author's first mention in the section, surname after;
%   (iii) pack first, but the rebuilt archon-cli-v3 index is an encouraged evidence surface —
%        used this revision for: self-citing author (citation-network §2), DA-03 roster (debate-map),
%        Quaranta claims (bor-sec-10), regulatory-paper theses (bor-sec-25/26), P1/P2 citations
%        (king-salvo-physio-2023 + grp-program-lineage), Freeman/Withy/McManus first names.
%
% BATCH 2 = ¶¶4–5 (drafted 2026-08-06, awaiting author review).
%   ¶4 (education): Mansikka three-way typology 262–66, second form + das Man 264 (bor-sec-18);
%     Gibbs datable-time/concealment 601–02 (bor-sec-07); Feldges & Pieczenko + Mertel +
%     Aroles & Küpers compactly in a footnote (bor-sec-16 and units); THOMSON GUARD APPLIED
%     per A-02 — cited for paideia/enframing/letting-learn frame, stated NOT a boredom source.
%   ¶5 (EEG): Kim/Seo/Laine null at 13 participants (A-06, bor-sec-28) — framed as the study
%     that sought agreement and documented disagreement; Barry et al. baselines 2765–66,
%     2771–72 (A-07, bor-sec-27); Yuvaraj et al. «Q16» at printed p. 11 (A-08, bor-sec-33)
%     with the no-resting-baseline gap; Miyauchi & Kawasaki, Perone, Yakobi, Seo compactly.
%   Names verified from index citations: Jan-Erik Mansikka, Paul Gibbs, Tom Feldges + Sonia
%     Pieczenko, Kevin Mertel [FIRST NAME ****** — index has "Mertel, K." only], Jeremy Aroles,
%     Wendelin Küpers, Iain Thomson, Joochan Kim, Jungryul Seo, Teemu H. Laine, Robert J. Barry,
%     Rajamanickam Yuvaraj, Eri Miyauchi, Masahiro Kawasaki, Sammy Perone, Ofir Yakobi.
%   Q16 is the batch's only rendered quotation. A-10 respected: no Miyauchi quotation exists.
%   Kim's null PARAPHRASED from the unit digest (no gated verbatim in bank) — E05 CONFIRMED.
%
% BATCH 1 = ¶¶1–3. REV 1 (2026-08-06, per author notes):
%   ¶1: Elpidorou NAMED as the self-citing author (citation-network.md §2: 7 of 8 A→C links are
%       Elpidorou citing Elpidorou; 8th = Hughes's critical footnote) + fn with both regulatory works
%       summarized; first "we" now names Christine King + fn listing P1/P2 and the appendix; ending
%       rebuilt around the Fahlman→Eastwood adoption and the design it governed.
%   ¶2: ****** resolved into DA-03 footnote (O'Brien/Hadjioannou/McManus + SZ 310 vs Withy + FCM
%       143/165; E&F periodizing at Affectivity I 662n2); Slaby point expanded via the disposedness
%       (Befindlichkeit) argument (link to Part I rendered through shared vocabulary, not explicit
%       back-reference; BT 172–75 locus carried from Part I draft); first-mention names applied.
%   ¶3: "VR headset"; conclusion rebuilt on Quaranta's actual apparatus (Loht's viewing-Dasein,
%       Drive My Car dead-time aesthetics, mood as neither work-alone nor viewer-alone, 31–32, 38).
%
% Batch-1 verification ledger:
%   Q14 printed p. 117 PINNED (folio-read: running head "The Other Side of Existence 117"; Slaby
%     fn 18 = Pessoa, Book of Disquiet, fragment 219). Q07 vol. pp. 207–08; Q12 printed p. 101.
%   P1 full citation from king-salvo-physio-2023-00-overview (ASEE 2023 #37129, 5 authors).
%   P2 from VLE manifest: "King et al., Assessment of Student Engagement in VR Clinical Immersion
%     Environments through Eye Tracking, ASEE 2024 #44685" — FULL AUTHOR LIST PENDING at assembly.
%   Elpidorou regulatory theses from bor-sec-25/26 digests (PhenomCogSci 17: 455–84; PhilPsych
%     31.3: 323–51). DA-03 roster from debate-map.md. Quaranta from bor-sec-10 (Film-Philosophy
%     28.1: 31–46); "neither subjective nor objective" p. 38 PARAPHRASED (quote not gated).
% PENDING VERIFICATION (assembly): first names for O'Brien, Hadjioannou (ed./ch. author), Hughes;
%   McManus chapter title as printed; Slaby 2010 venue for WC; SZ-vs-BT page convention harmonize.
% OPEN: Appendix designation for P1/P2 reprints (****** in fn A). E01 open item CLOSED this rev
%   (author = Elpidorou, index-verified).
% Selection decisions carried: Q01 barred; Q08 held for ¶10 (paraphrased here); Q09 paraphrased;
%   Q10 unused (optional); Q11 skipped; Q13 avoided (hazard); O-06 = cite-not-re-present.
%
% WORKS CITED introduced by this batch (for assembly; ****** = verify there):
%   Eastwood, John D., et al. "The Unengaged Mind: Defining Boredom in Terms of Attention."
%     Perspectives on Psychological Science, 2012. [vol/no/pp ****** at assembly]
%   Elpidorou, Andreas. "The Bored Mind Is a Guiding Mind: Toward a Regulatory Theory of Boredom."
%     Phenomenology and the Cognitive Sciences, vol. 17, 2018, pp. 455-84.
%   ---. "The Good of Boredom." Philosophical Psychology, vol. 31, no. 3, 2018, pp. 323-51.
%   Elpidorou, Andreas, and Lauren Freeman. "Affectivity in Heidegger I: Moods and Emotions in
%     Being and Time." Philosophy Compass, 2015. [vol/pp ******]
%   ---. "Is Profound Boredom Boredom?" Heidegger on Affect, ed. Hadjioannou, 2019, pp. 207-[**].
%   Fahlman, Shelley, et al. "Development and Validation of the Multidimensional State Boredom
%     Scale." 2013. [journal details ******]
%   Hughes [first name ******]. "Meaninglessness and Monotony in Pandemic Boredom." 2023.
%   King, Christine E., Kit Roy Feeney, Quangminh Tang, Milan Das, and Dalton Salvo. "Work in
%     Progress: Physiological Assessment of Learning in a Virtual Reality Clinical Immersion
%     Environment." ASEE, 2023. Paper #37129.
%   King, Christine E., et al. [****** full list] "Assessment of Student Engagement in VR Clinical
%     Immersion Environments through Eye Tracking." ASEE, 2024. Paper #44685.
%   O'Brien [******]. "Being, Nothingness and Anxiety." Heidegger on Affect, 2019.
%   Hadjioannou [******], ed. Heidegger on Affect. Rowman?, 2019. [publisher ******] — also ch. 4,
%     "Angst as Evidence: Shifting Phenomenology's Measure."
%   McManus, Denis. [ch. 6 title per index: "Affect and Authenticity"] Heidegger on Affect, 2019.
%   Withy, Katherine. "Finding Oneself, Called." Heidegger on Affect, 2019.
%   Quaranta, Chiara. "In the Mood for Heideggerian Boredom? Film Viewership as
%     Being-in-the-World." Film-Philosophy, vol. 28, no. 1, 2024, pp. 31-46.
%   Slaby, Jan. "The Other Side of Existence: Heidegger on Boredom." 2010, pp. 101-20. [venue ******]
% --- introduced by batch 2 ---
%   Aroles, Jeremy, and Wendelin Küpers. "Towards an Integral Pedagogy in the Age of 'Digital
%     Gestell'." 2022. [venue/pp ******]
%   Barry, Robert J., Adam R. Clarke, Stuart J. Johnstone, Christopher A. Magee, and Jacqueline
%     A. Rushby. "EEG Differences Between Eyes-Closed and Eyes-Open Resting Conditions."
%     Clinical Neurophysiology, vol. 118, 2007, pp. 2765-73.
%   Feldges, Tom, and Sonia Pieczenko. "Boredom in Educational Contexts: A Critical Review."
%     Encyclopaideia, vol. 24, no. 57, 2020, pp. 1-15.
%   Gibbs, Paul. "The Concept of Profound Boredom: Learning from Moments of Vision." Studies in
%     Philosophy and Education, vol. 30, 2011, pp. 601-13.
%   Kim, Joochan, Jungryul Seo, and Teemu H. Laine. "Detecting Boredom from Eye Gaze and EEG."
%     Biomedical Signal Processing and Control, vol. 46, 2018, pp. 302-13.
%   Mansikka, Jan-Erik. "Can Boredom Educate Us? Tracing a Mood in Heidegger's Fundamental
%     Ontology from an Educational Point of View." Studies in Philosophy and Education, vol. 28,
%     2009, pp. 255-68.  [2008/2009 — confirm year at assembly ******]
%   Mertel, K[****** first name]. "Heidegger, Technology and Education." Journal of Philosophy
%     of Education, vol. 54, no. 2, 2020, pp. 467-86.
%   Miyauchi, Eri, and Masahiro Kawasaki. "Electroencephalogram Oscillations Support the
%     Involvement of Task-Unrelated Thoughts in the Mechanism of Boredom: A Pilot Study." 2018.
%     [venue/pp ******]
%   Perone, Sammy, Elizabeth H. Weybright, and Alana J. Anderson. "Over and Over Again: Changes
%     in Frontal EEG Asymmetry Across a Boring Task." 2019. [venue/pp ******]
%   Seo, Jungryul, Teemu H. Laine, and Kyung-Ah Sohn. "Machine Learning Approaches for Boredom
%     Classification Using EEG." 2019. [venue/pp ******]
%   Thomson, Iain. "Heidegger on Ontological Education, or: How We Become What We Are."
%     Inquiry, vol. 44, 2001. [pp ******]
%   Yakobi, Ofir, James Boylan, and James Danckert. "Behavioral and Electroencephalographic
%     Evidence for Reduced Attentional Control and Performance Monitoring in Boredom." 2021.
%     [venue/pp ******]
%   Yuvaraj, Rajamanickam, et al. "EEG-Based Functional Connectivity Patterns During Boredom in
%     an Educational Context." Scientific Reports, vol. 15, 2025, article 35439.

Boredom has been academically pursued across two rather different fields of study. One asks
what boredom discloses about the person who suffers it, treats the experience as a disclosure
of how things stand for someone rather than as a state to be relieved, and reaches for Martin
Heidegger and the phenomenological tradition that follows him; the other asks what boredom
looks like in a signal, and reaches for electrodes, eye trackers, and self-report scales
validated against one another. Across the works surveyed here, the connection between them
amounts to one philosopher, Andreas Elpidorou, citing his own work across his own two research
programs, together with one critical footnote written by someone else.\footnote{Elpidorou
publishes in both literatures, and the phenomenological essays---several written with Lauren
Freeman---cite two papers of Elpidorou's own regulatory-theory program. ``The Bored Mind Is a
Guiding Mind: Toward a Regulatory Theory of Boredom'' (\textit{Phenomenology and the Cognitive
Sciences}, 2018) argues that boredom is a functional emotion, at once informing a person that
the present situation has stopped satisfying and motivating the pursuit of a new goal. ``The
Good of Boredom'' (\textit{Philosophical Psychology}, 2018) argues that transitory state
boredom serves well-being precisely by keeping a person in motion toward goals worth having.
The one bridge not of Elpidorou's making runs the opposite direction and is a critique:
Hughes, ``Meaninglessness and Monotony in Pandemic Boredom'' (2023), objects in a footnote
that Elpidorou's normative-psychological turn departs from Heidegger's project.} The first
literature counts a reading of experience
as evidence; the second counts a number extracted from a sensor. Neither treats what the
other counts as bearing on its own question, and the absence is not an oversight either could
have corrected by wider reading. The two studies Christine King and I published stand in
neither literature:\footnote{The two studies: King, Christine E., Kit Roy Feeney, Quangminh
Tang, Milan Das, and Dalton Salvo, ``Work in Progress: Physiological Assessment of Learning
in a Virtual Reality Clinical Immersion Environment'' (ASEE 2023, Paper \#37129), the
three-participant pilot that first instrumented the platform's videos---EEG, cardiac, and
gaze---against boring and interesting control films; and our eye-tracking study of student
engagement on the same platform, ``Assessment of Student Engagement in VR Clinical Immersion
Environments through Eye Tracking'' (ASEE 2024, Paper \#44685). Both papers are reproduced in
full in Appendix ******.} they belong to a third lineage, clinical immersion in engineering
education, where a headset substitutes for clinical access that no program can secure in
sufficient quantity, and that lineage arrives with its own question, its own warrant, its own
reference list, and no commitment about what boredom is. From the second literature we took
one thing, deliberately and by name: the definition of boredom under which we collected the
data. Our first paper adopts Shelley Fahlman's formulation of it---the definition
underwriting the Multidimensional State Boredom Scale, which carried John Eastwood's
attentional account of boredom into measurement---and the adoption governed the design end to
end. Because the definition fixes boredom as an aversive episode, an unfulfilled desire for
engagement that would satisfy, the experiment treats boredom as something a film can induce
and relieve: it sets a film selected to starve engagement beside a film selected to feed it
and beside the clinical footage the platform exists to deliver, and it points its
instruments---EEG, eye tracking, and a per-episode self-report---at what each film does to
the person watching. What the definition does not contemplate is a boredom that is not about
the film at all. The data collected under it contains exactly that.

% ****** AUTHOR ACTION PENDING (¶1 footnote b) — rendered so it cannot be lost in a compile.
% ****** Remove this box once the appendix exists and the footnote's ****** is filled.
\begin{center}
\fbox{\parbox{0.9\linewidth}{\centering\textbf{****** AUTHOR ACTION PENDING ******}\\
Add the appendix containing the two published articles (P1, ASEE 2023 \#37129; P2, ASEE 2024
\#44685), then fill the appendix designation in the footnote above. Also pending: P2's full
author list for the works cited.}}
\end{center}

The phenomenological literature's liveliest quarrel concerns whether the experience Heidegger
analyzes is boredom at all. Elpidorou and Lauren Freeman set profound boredom against the
construct the empirical programs measure---a brief, situation-bound, many-sided response to
unsatisfying circumstances---and conclude of profound boredom that ``«Q07»'' (207--08). The
verdict protects both parties to the comparison: it spares the psychological construct a
metaphysics it never claimed, and it cautions every instrumented study, this one included,
against announcing that a sensor has caught what Heidegger meant. Nor is it settled which
mood holds the ground position in Heidegger's own architecture: the case for anxiety
(\textit{Angst}) as the fundamental attunement (\textit{Grundstimmung}) rests on
\textit{Being and Time}'s near-unique designation, the case for profound boredom on
\textit{The Fundamental Concepts of Metaphysics}, and Elpidorou and Freeman observe that the
question is partly one of period, since the term appears only once in \textit{Being and Time}
and becomes load-bearing only after 1929.\footnote{The anxiety side is argued across three
chapters of \textit{Heidegger on Affect} (ed. Hadjioannou, 2019): O'Brien, ``Being,
Nothingness and Anxiety''; Hadjioannou, ``Angst as Evidence: Shifting Phenomenology's
Measure''; and Denis McManus on affect and authenticity---with \textit{Being and Time}
itself giving anxiety the designation nearly alone (SZ 310). Katherine Withy's chapter in the
same volume, ``Finding Oneself, Called,'' makes the affirmative case that profound boredom
can occupy the role, with support at \textit{FCM} 143 and 165. For the periodizing
observation: Elpidorou and Freeman, ``Affectivity in Heidegger I: Moods and Emotions in
\textit{Being and Time}'' (2015), 662n2.} Jan Slaby reads the whole affective register as
``«Q12»'' (101): mood is not noise laid over thinking but the way a person's standing in the
world shows itself at all. Slaby's formula names what \textit{Being and Time} calls
disposedness (\textit{Befindlichkeit})---one does not first register a neutral scene and
afterward color it with feeling; one finds oneself already attuned, affected by what one
encounters before any reflective act of judgment arrives (\textit{BT} 172--75)---and the
formula carries a consequence. If attunement is disclosure rather than coloration, then a
person can remain occupied---attentive, compliant, even performing well---while what the
attunement discloses is emptiness. Fernando Pessoa's words, as Slaby borrows them, give that
depth its starkest formula---``«Q14»'' (qtd. in Slaby 117). An occupied surface over a hollow
depth is, under this reading, a structure open to description, not a contradiction in terms.

Film is the point where the phenomenological literature reaches this experiment's own medium.
Chiara Quaranta, building on Shawn Loht's account of film viewing as a genuine mode of
being-in-the-world---a two-way disclosure in which the viewer opens onto the film and the
film discloses to the viewer---argues that Heidegger's three forms of boredom, the profound
form above all, name what certain films do to the people watching them. The case study is
Ryusuke Hamaguchi's \textit{Drive My Car}, whose fixed shots, silences, repetitions, and dead
time Quaranta reads as an aesthetics composed to arouse profound boredom---with the standing
qualification that no film accomplishes this alone, since a mood arises neither from the work
by itself nor from the viewer by itself but between the two (31--32, 38). Within the works
surveyed here this is the one sustained reading of film through the boredom analysis, and the
arrangement it describes is the arrangement this experiment builds and instruments: a film
selected for monotony, watched under mandate rather than by choice, in a VR headset that
holds the footage before the eyes, while EEG, eye tracking, and a per-episode self-report
record what the arrangement does to the person inside it. Quaranta reads the arrangement off
a finished film; the laboratory assembles it on purpose and asks every channel to report.

Education is where this analysis has been put to work longest, and where its most useful
precedent for the present study is found. Jan-Erik Mansikka fits Heidegger's three forms to
three ways of learning: the first to technical instruction that reduces phenomena to
abstractions and leaves the student confronting dull matter; the second to a learning that is
meaningful and still leaves the student empty; the third to boredom as a possible rupture,
which Mansikka is careful not to prescribe as method (262--66). The second is the one that
matters here. A student, on his account, can be outwardly exemplary---fully occupied,
participating, even high-achieving---while what governs the performance is the anonymous
``they'' (\textit{das Man}) rather than any self-relation to the material (264). The
diagnostic consequence follows directly: engagement and this hollow boredom are compatible
rather than opposed, so a metric that counts engagement---time on task, submissions,
participation---cannot detect the second form, because everything such a metric registers is
precisely what the second form leaves intact. Mansikka reaches that conclusion in a classroom
register more than a decade before this experiment; what the laboratory adds is two channels
that can disagree, and the disagreement in the data is that argument in instrumented form.
Paul Gibbs pursues the same third form toward pedagogy, reading modern datable time---the
semester, the credit hour, the examination date---as the very thing that conceals boredom's
disclosive work, so that busyness is not boredom's opposite but the mechanism of its
concealment (601--02). Tom Feldges and Sonia Pieczenko press a different objection, that the
field's dominant taxonomy of academic boredom conflates what precedes boredom with boredom
itself, and that classroom boredom may be an artifact of frontal instruction rather than
anything intrinsic to learning; Kevin Mertel, Jeremy Aroles, and Wendelin Küpers extend the
enframing critique to educational technology.\footnote{Feldges and Pieczenko, ``Boredom in
Educational Contexts: A Critical Review,'' \textit{Encyclopaideia} 24.57 (2020); Mertel,
``Heidegger, Technology and Education,'' \textit{Journal of Philosophy of Education} 54.2
(2020); Aroles and K\"upers, ``Towards an Integral Pedagogy in the Age of `Digital
\textit{Gestell},''' 2022.} One frequent citation in this neighborhood requires a
correction. Iain Thomson's account of ontological education is standardly grouped with these
sources, and it supplies the frame within which they argue---\textit{paideia} against
enframing, teaching as letting learn, and the university's crisis as ontohistorical rather
than institutional---but the words ``bored,'' ``boredom,'' and \textit{Langeweile} do not
occur anywhere in his text. It is cited here for the frame it supplies and not as a reading
of boredom.

The instrumented literature arrives at this experiment's own channels, and arrives carrying
its nulls. Joochan Kim, Jungryul Seo, and Teemu Laine ran the closest precursor to this
study's design: a monotonous video, a consumer EEG headband and a screen-based eye tracker
recording together, and a post-video boredom rating. Their thesis was that the two
instruments would agree with each other during boredom. Temporally coincident regions did
appear in both signals across all thirteen analyzed participants, and the self-report
confirmed that boredom had been elicited---but that rating correlated significantly with no
extracted physiological feature at all. The study that set out to demonstrate cross-channel
agreement instead documented the disagreement, which is the finding this section inherits and
extends. Robert Barry and colleagues supply the methodological floor beneath every EEG claim
made here: eyes-closed and eyes-open rest are not interchangeable baselines, since alpha power
indexes global arousal while the focal bands index activation, so a study must choose its
baseline and defend the choice (2765--66, 2771--72). That bears on this experiment twice, once
on the EEG negative and once on eye closure as a surface in its own right. Rajamanickam
Yuvaraj and colleagues are the nearest neighbor in both stimulus and setting---university
students, a boring lecture video, whole-scalp connectivity---and their finding is that boredom
is not a network-idle state: gamma-band connectivity grew more locally clustered and less
globally efficient, which they read as ``«Q16»'' rather than mere disengagement (11). Their
design also carries the gap this section must not repeat, since both of their conditions are
active eyes-open video-viewing with no resting baseline of any kind. Eri Miyauchi and Masahiro
Kawasaki on task-unrelated thought, Sammy Perone and colleagues on frontal asymmetry across a
boring task, Ofir Yakobi and colleagues on attentional control, and Jungryul Seo and colleagues
on machine classification fill out the strand.
